% Bioinformatics Report Template — Differential Expression with DESeq2
% Extends the core mycelium report template with bioinformatics-specific sections.

\documentclass[11pt,a4paper]{article}

\usepackage[utf8]{inputenc}
\usepackage[T1]{fontenc}
\usepackage{amsmath,amssymb}
\usepackage{graphicx}
\usepackage{booktabs}
\usepackage{hyperref}
\usepackage[margin=1in]{geometry}
\usepackage{natbib}
\usepackage{xcolor}
\usepackage{caption}
\usepackage{subcaption}
\usepackage{longtable}         % For gene tables spanning multiple pages

\title{Differential Expression Analysis: [Title]}
\author{[Author(s)]}
\date{\today}

\begin{document}

\maketitle

\begin{abstract}
[Summarize: biological question, experimental design (N samples, conditions),
key DE results (N genes at FDR threshold), top enriched pathways, main conclusion.]
\end{abstract}

\section{Introduction}

[Biological context. Why was this analysis performed? What is the hypothesis
or question? Reference relevant prior work.]

\section{Methods}

\subsection{Experimental Design}

[Describe samples, conditions, replicates, known covariates. Include a table
if helpful.]

% \begin{table}[htbp]
%     \centering
%     \caption{Sample overview.}
%     \begin{tabular}{llll}
%         \toprule
%         Sample ID & Condition & Batch & Notes \\
%         \midrule
%         S1 & treated & B1 & \\
%         \bottomrule
%     \end{tabular}
% \end{table}

\subsection{Data Processing}

\begin{itemize}
    \item \textbf{Genome}: [assembly, annotation version]
    \item \textbf{Alignment}: [STAR/HISAT2/Salmon version and parameters]
    \item \textbf{Quantification}: [featureCounts/tximport version]
    \item \textbf{QC}: FastQC v[X], MultiQC v[X]
\end{itemize}

\subsection{Differential Expression}

DESeq2 v[version] \citep{Love2014} was used for differential expression analysis.
Genes with fewer than [10] total counts across all samples were filtered before
analysis. The design formula was \texttt{$\sim$ [covariates + condition]}.
Log fold change shrinkage was applied using the apeglm method \citep{Zhu2019}.
Genes were considered differentially expressed at an adjusted $p$-value (Benjamini-Hochberg)
$< 0.05$ and $|\log_2\text{FC}| > 1$.

\subsection{Enrichment Analysis}

[Method: fgsea/clusterProfiler. Gene sets: MSigDB Hallmark/KEGG/Reactome.
Describe ranking metric and significance threshold.]

\section{Results}

\subsection{Quality Control}

[Summarize QC results. Reference MultiQC report in supplementary. Note any
samples excluded and why.]

% \begin{figure}[htbp]
%     \centering
%     \includegraphics[width=0.6\textwidth]{../outputs/figures/pca_plot.pdf}
%     \caption{PCA of normalized expression values, colored by condition.}
%     \label{fig:pca}
% \end{figure}

\subsection{Differential Expression}

[Report number of DE genes. Show volcano plot and MA plot.]

% \begin{figure}[htbp]
%     \centering
%     \begin{subfigure}{0.48\textwidth}
%         \includegraphics[width=\textwidth]{../outputs/figures/volcano_plot.pdf}
%         \caption{Volcano plot.}
%     \end{subfigure}
%     \begin{subfigure}{0.48\textwidth}
%         \includegraphics[width=\textwidth]{../outputs/figures/ma_plot.pdf}
%         \caption{MA plot.}
%     \end{subfigure}
%     \caption{Differential expression results. Significant genes (FDR $< 0.05$,
%     $|\log_2\text{FC}| > 1$) highlighted in red.}
%     \label{fig:de}
% \end{figure}

% Top DE genes table
% \begin{table}[htbp]
%     \centering
%     \caption{Top 10 differentially expressed genes by adjusted $p$-value.}
%     \begin{tabular}{lrrr}
%         \toprule
%         Gene & $\log_2$FC & $p_{\text{adj}}$ & Mean Expression \\
%         \midrule
%         [gene] & & & \\
%         \bottomrule
%     \end{tabular}
% \end{table}

\subsection{Pathway Enrichment}

[Report top enriched pathways. Show enrichment dot plot.]

% \begin{figure}[htbp]
%     \centering
%     \includegraphics[width=0.8\textwidth]{../outputs/figures/enrichment_dotplot.pdf}
%     \caption{Top enriched pathways by normalized enrichment score.}
%     \label{fig:enrichment}
% \end{figure}

\section{Discussion}

[Interpret results in biological context. Compare with prior work.
Discuss limitations (sample size, batch effects, genome annotation).
Propose next steps.]

\bibliographystyle{plainnat}
\bibliography{references}

\end{document}
